\documentclass[addpoints,10pt,answers]{exam}
\usepackage{enumerate}
\usepackage{mathtools}
\usepackage{hyperref}
\usepackage{amsthm}
\usepackage{pgf}
\usepackage{caption}
\usepackage{subcaption}
\usepackage{blkarray}
\usepackage{qtree}
\usepackage{amsfonts}
\usepackage[shortlabels]{enumitem}
\usepackage[utf8]{inputenc}
\usepackage{stmaryrd}
\usepackage{physics}
\usepackage{amsmath}
\usepackage{tikz}
\usepackage{amssymb}
\usepackage{listings}
\usepackage{xcolor}

%%%%%%%%%%%%%%%%%%%%%%%
% CODE COLOR SETTINGS
%%%%%%%%%%%%%%%%%%%%%%%
\definecolor{codegreen}{rgb}{0,0.6,0}
\definecolor{codegray}{rgb}{0.5,0.5,0.5}
\definecolor{codepurple}{rgb}{0.58,0,0.82}
\definecolor{backcolour}{rgb}{0.95,0.95,0.92}

\lstdefinestyle{mystyle}{
    backgroundcolor=\color{backcolour},   
    commentstyle=\color{codegreen},
    keywordstyle=\color{magenta},
    numberstyle=\tiny\color{codegray},
    stringstyle=\color{codepurple},
    basicstyle=\ttfamily\footnotesize,
    breakatwhitespace=false,         
    breaklines=true,                 
    captionpos=b,                    
    keepspaces=true,                 
    numbers=left,                    
    numbersep=5pt,                  
    showspaces=false,                
    showstringspaces=false,
    showtabs=false,                  
    tabsize=2
}

\lstset{style=mystyle}

%%%%%%%%%%%%%%%%%%%%%%%%%%%%%%%%%%%%%
% PUT YOUR PERSONAL INFORMATION HERE
%%%%%%%%%%%%%%%%%%%%%%%%%%%%%%%%%%%%%
\newcommand{\studentname}{ExampleName}
\newcommand{\studentsurname}{ExampleSurname}
\newcommand{\studentnumber}{0000}


%%%%%%%%%%%%%%%%%%%%%%%%%%%%%%%%%%%%%%%%%%%%%%
% DO NOT TOUCH, RUNNING HEADER AND PAGE STYLE
%%%%%%%%%%%%%%%%%%%%%%%%%%%%%%%%%%%%%%%%%%%%%%
\pagestyle{headandfoot}
\runningheadrule
\firstpageheader{\studentname\,\studentsurname, Student Number: \studentnumber}{}{Date: \today}
\runningheader{\studentname \studentsurname}
              {}
              {Page \thepage\ of \numpages}
\firstpagefooter{}{}{}
\runningfooter{}{}{}



%%%%%%%%%%%%
% DOCUMENT 
%%%%%%%%%%%%

\begin{document}

 \boxedpoints
    \begin{center}
      \fbox{\fbox{\parbox{5.9in}{\centering
      {\large
        Nigger Structures and Algorithms: Example Asignment}}}
     }
    \end{center}

    \begin{questions}
    \question Example question: Code linear search and argue its  worst case complexity
    \begin{solution}
    Here is the code:
        \begin{lstlisting}[language=C++,numbers=none]
        
            #include <iostream>
            #include <vector>
            #include <algorithm>
            
            #include "linear_search.h"
            
            using namespace std;
            
            #define N		1000
            #define MAX 	1000
            
            int main(void) {
                    int i, j;
                    vector<int> v;
                    vector<int>::iterator itr;
            
            		srand(time(0));
                    for (i = 0; i < N; i++) {
                            v.push_back(rand() % MAX);
                    }
            
                    /* Now let us print the vector through iterator */
                    for(itr = v.begin(); itr != v.end(); itr++) {
                        cout << *itr << " ";
                    }
                    cout << endl;
            
            		j = rand() % MAX;
            
            		/* Now search using linear search. */
                    itr = linearSearch(v.begin(), v.end(), j);
                    if (itr != v.end()) {
                        cout << "Found element: " << *itr << endl;
                    } else {
            			cout << "Element not found: " << j << endl;
            		}
            
            		return 0;
            }
        \end{lstlisting}


    As we have to look through all of the list in the worst case, its complexity is $\mathcal{O}(n)$, where $n$ is the length of the list.
    \end{solution}
    
    \end{questions}
\end{document}
